% =============================================================================
% warn-overfull-line -- an honest mistake: one unbreakable \texttt{} string
% overflows slightly into the margin (a classic overfull \hbox).
%
% The overflow is a few points, far from deliberate squeezing, so the
% checker reports it for a human glance without failing the paper.
%
% Expected checker outcome: PASS, with warning [text-block].
% =============================================================================
\documentclass[letterpaper,twocolumn,10pt]{article}
\usepackage{usenix}
\usepackage{lipsum}

\date{}
\title{\Large \bf A Submission with One Overfull Line}
\author{{\rm Alice Anonymous}\\One University \and {\rm Bob Blinded}\\Two Institute}

\begin{document}
\maketitle

\begin{abstract}
\lipsum[1][1-5]
\end{abstract}

\section{Introduction}
\lipsum[2-4]

\begin{figure*}[t]
  \noindent
  \mbox{\texttt{a-very-long-unbreakable-command-line-option-string-that-does-not-quite-fit-into-the-text-block-ok-yes}}
  \caption{An unbreakable command line that pokes into the margin.}
\end{figure*}

\section{Design}
\lipsum[5-7]

\appendix
\section*{Ethical Considerations}
This synthetic test paper raises no ethical concerns.

\section*{Open Science}
All artifacts of this test live in the checker repository.

\begin{thebibliography}{1}
\bibitem{esx} C.~A.\ Waldspurger. Memory resource management in VMware ESX
  server. In \emph{OSDI}, 2002.
\end{thebibliography}
\end{document}
